% !TeX root = ../papers/nc-p3.tex
\section{三次の可約な階数一配置の全分割}
\label{sec:rank-one-cubic-arrangement-exclusion}
\begin{theorem}\label{thm:rank-one-cubic-arrangement-exclusion}
被約な四辺有限の標準有限Frobenius窓で、一般可逆な全階数反対根の階数一軌跡の被約曲線部分は次数2以下である。
\end{theorem}
\begin{proof}[連結三次の三つの二次式]
三次の連結和 $C$ は平面内に入らず、平面二次部分 $D$（交わる二直線でもよい）とそこに会う直線 $J$ の和である。
前定理の木による切断空間は四次元なので二つの座標四空間は平方同型とも両立して同一となる。
よって合同基底で $N=xx^t+A(q)$。$D$ の平面を $h=0$、$J$ の二式を $l,m$ として
$q=(hl,hm,Q)$ は $C$ のイデアルを生成する。実際 $Q|_{h=0}$ は $D$ の式で、
$D$ と $J$ の交点で消えるから $Q\in(l,m)$ となるよう延長できる。$(h,Q)\cap(l,m)=(hl,hm,Q)$ である。

六つの交代評価を張るには $A$ は四直線判別式をもつ対角網となる。
定理\ref{thm:rank-one-twisted-cubic-exclusion}の証明で、$q$ の繊維によらないことを確認する。
一般階数3の定核は列の平面因子で排除され、動く原始核は一次である。
このとき階数二パラメータは直線と高々一点に入る。直線上で一般階数2なら第二外冪の原始Plücker列は次数1以下なので
方向の包は高々二次元であり、一般階数1以下のときは寄与しない。
よって一点を加えても六方向には足りない。一般階数4での四次判別式の分割と対角化は同定理の証明そのものである。
従って各座標辺 $E_{ij}$ が指定された値 $t_{ij}$ に対する曲線外印を一つ以上持ち、
これらは全ての曲線外印を覆う。どの辺も $C$ の成分にならない。

$E_{ij}$ 上に曲線外印が二つ以上あれば、$q$ のこの辺での値は一定 $t_{ij}$ であり、$C$ と長さ二で交わる。
実際 $C$ に素な辺の場合、$C$ とその直線を同時に含む二次曲面はない（前定理最後の連結三次と素な直線の議論）。
従って三本の $q$ はその直線で全二次式を与え単射な写像となる。切断長1なら固定因子を除いて一次の完備系、
切断長2なら一定となる。長さ3以上は辺全体が $C$ に入るので不可能である。
値一定の辺には二つの曲線外印までしか載らない。二行小行列式が辺上の四次式で、既に長さ二で $C$ を含み、
恒等零なら別の階数一直線が必要になるからである。
$C$ 上は合計九印までなので曲線外は十一印以上となり、六辺中五本以上で値が一定となる。
各座標頂点には異なる二つの定値辺が集まり、よって $q$ がそこで定義されず頂点は $C$ 上にある。
すると全六辺が頂点二つを含み $q$ は各辺で指定の一定値となる。各頂点での二つの零点は単純である。
従って一点 $e_i$ からの三辺の方向微分 $dq(e_i)$ の値はそれぞれ $t_{ij}$ の非零倍となる。
しかし曲線 $C$ 上のJacobi行列の階数は高々2で、これら三つのパラメータは判別四直線の残る三直線の三角形で非共線である。
これで連結三次の場合を除いた。
\end{proof}
\begin{proof}[平面二次と素な直線]
次に $D\cap J=\varnothing$ とする。前定理の $k$ の列では $D$ 上で消える四二次式のうち三つが $h$ の倍数で、
適切な独立三つ組により $k m_i(y)$ が $h l_i(x)$ に対応する。
入力平面と $J$ の交点を $r$、出力平面と $J'$ の交点の逆像を $s$ とすると $r\ne s$ である。
実際 $h=k=0$ の階数一集合は入力平面上の $q Bk=0$ から $D$ のみである。
$J$ 上の $s$ が $l_i$ を全て消し、$r'$（$r$ の出力）が $m_i$ を全て消す。
従って各三つ組はそれぞれ $s,r'$ を中心とする平面への射影である。
平面二次と直線上でそれぞれ座標を合わせた共通切断も三つある。
実際前定理の次元評価で $d+b=3+2=5$ である。$D$ を通る二次式が $J$ でも零なら、
$h=0$ に制限した $D$ の式の係数は $r\notin D$ により零となる。
よってその全空間は $h(l,m)$ で、$l,m$ は $J$ の二つの一次式である
（この $l,m$ は先の三つ組とは別である）。
三つの共通切断での平方の差は対称三行鉛筆 $hA(l,m)$ となる。
鉛筆が特異なら固定核または零二行ブロックによる
定理\ref{thm:rank-one-quartic-normalization}の矛盾が使える。
一般可逆なら $h=0$ と $J$ の外の全印は三つ以下の $J$ を通る平面に入る。

各々の平面 $P$ について、そこに属す $h=0,J$ 外の印の出力も $J'$ を通る一平面 $P'$ に入る。
$k=0$ 上の $D$ 外の階数一点は高々一つで、それは既に $s$ の出力であるから、これらの印では $h,k$ ともに非零である。
従って上の二つの三つ組による中心射影は同じ値を与える。$P$ は $s$ を含んで直線に射影され、
対応する出力平面は $r'$ を含む。また $J,J'$ は各中心射影で同じ一点に潰れるためこの平面は $J'$ を含む。
二平面 $P,P'$ の二次の包は一致できない（一致すれば $P'$ の零化列が $P$ で消える）。
従って同時にこの二平面に属する入力データは $P$ 内の非零二次式で消える。
それは $J$ の全体で消えるため、平面 $P$ の該当印は $J$ の外で一直線上の三つ以下である。
入力平面 $h=0$ の $D$ 外の階数一点は既に $r\in J$ のみなので、
全印数は $D$ 上の六つ、$J$ 上の三つ、三平面で各三つの合計十八以下となって矛盾する。
\end{proof}
\begin{proof}[互いに素な三直線の二つずつの射影]
残る三直線 $J_1,J_2,J_3$ が互いに素とする。同じ母線族の唯一の二次曲面 $Q$ を
$\mathbb P^1\times\mathbb P^1$ とし各 $J_i$ では同じ第二パラメータ $u$ を使う。
出力でも直線は互いに素である。二次の包 $\operatorname{Sym}^2 J_i$ 同士のベクトル空間の交わりが零かどうかで
直線の交わりが分かるからである。対応する $Q'$ を同様に取り、各 $J_i$ での出力の第二パラメータを $A_i u$ と書く
（$A_i$ は二変数の射影同型）。$Q$ 上で三直線外の階数一点は高々二つである。
実際 $u$ 固定の横断母線上では二次行列 $N$ が三点の階数一行列のLagrange補間である。三点以外で三係数は非零なので、
三出力ベクトルが独立なら階数3である。従属なら同じ横断出力母線上で $A_1u=A_2u=A_3u$ となる。
三射影同型が等しければ $Q$ 全体で階数二以下となるので不可能、従ってこの方向は高々二つ。
各母線では四次小行列式に三つの既知零点があり、全母線が階数一なら総次数上界に反するから、
さらに一つまでしかない。

対 $i,j$ に対し各直線を核とする二座標への射影を $x_i,x_j$、出力の同様の射影を $y_i,y_j$ とする。
相補な直線なので $(x_i,x_j)$ は入力の二組の座標になり、$x_i\otimes x_j$ の四成分は
$J_i,J_j$ 上で消える二次式全部を与える。よって二次座標同型による可逆な四行変換 $T_{ij}$ があり、
全階数一点で $y_i\otimes y_j$ は $T_{ij}(x_i\otimes x_j)$ に正規化を除き等しい。
各 $x_i$ は $Q$ 上でその零点を除き $u$ と同じ射影座標に取れ、$y_i$ も出力の第二パラメータで取る。
四次式 $\det T_{ij}(x_i\otimes x_j)$ は $Q$ 上で消える。実際 $x_i\otimes x_j$ はそこで
$u\otimes u$ の倍であり、第三の直線上で既に全ての $u$ に対する階数一が確認できる。
この式が非零なら $Q$ を除いた残余 $G_{ij}$ は二変数組で $(1,1)$ の二次式で $J_i,J_j$ を含み、
$Q$ 外の全印を含む。第三の直線 $J_k$ の二つの独立な零化一次式を $G_{ij}$ に掛けると、
$Q$ 上の三母線外の高々二点以外で全て消える三次式の二次元空間を得る。
三次評価の同型からその空間の上界は2である。異なる二対で $G_{ij}$ が非零なら両空間は等しく、
最大公約因子 $G_{ij}$ が比例し、それを除いた一次式鉛筆も同一になる。異なる第三直線に反する。

従って少なくとも二対、添字を選び $(1,2),(1,3)$ で $T_{ij}$ はSegre二次曲面を保つ。
この自己同型は二族の母線を保つか交換するかである。交換なら $y_i$ は $x_j$ の固定射影変更と比例し、
その二項交代式が全印で消える（$J_i$ 上で $y_i=0$、$J_j$ 上で $x_j=0$）。混合評価に反する。
族を保つ場合、第三直線上での値により、$i,j$ の外では $y_i$ は $A_k x_i$ に比例し $y_j$ は $A_k x_j$ に比例する。
この式の各交代評価は非零の双一次式であり、従って $J_i$ と $J_j$ はそれぞれ少なくとも一つの指定印を持つ。
$A_2=A_3$ なら $y_1$ と $A_2 x_1$ の交代式が二対の式で全印を覆って零となり矛盾する。
異なれば三母線外の全印は $A_2 x_1=A_3 x_1$（射影値）を満たす。従って $J_1$ を通る高々二平面に入る。
その二平面の和には $J_1$ の印の重複を引いて高々 $14-1=13$ 印しかない。
残る $J_2,J_3$ 上の各三印を加えても十九以下となる。この矛盾で最後の分割も除いた。
\end{proof}
